\section{Background \& Related Work}
\label{sec:background}
This section introduces the background on text-to-image diffusion models and reviews existing approaches for concept unlearning in diffusion models.

\subsection{Background}
% \varun{what is the significance of the discussion around diffusion mdoels? why is this primer needed to understand the rest of your paper? need to give that context!}
We briefly introduce the latent diffusion notation used throughout the paper, since our analysis of concept activation regions and unlearning-induced perturbations is defined over the model's latent, text-conditioned denoising process.

\textbf{Text-to-image latent diffusion models.}
We consider a text-to-image latent diffusion model (LDM) of the form introduced in Stable Diffusion~\citep{rombach2022high}.
% We consider a text-to-image latent diffusion model (LDM), following Stable Diffusion~\citep{rombach2022high} \varun{what does it mean to follow?}. 
An input image $\mathbf{x} \in \mathcal{X}$ is first mapped by a pretrained variational autoencoder encoder $E$ into a latent representation 
% \varun{the latent z is bold-faced in the next line, not bold-faced here? subsequent usages are also mixed in usage} 
$\mathbf{z} = E(x) \in \mathcal{Z}$, and reconstructed by a decoder $V : \mathcal{Z} \rightarrow \mathcal{X}$. 

% \ali{I think this is extra info: 
% Instead of modeling pixels directly, the diffusion model operates in latent space, which substantially reduces computation while preserving semantic content.}
Given a latent $\mathbf{z}$, the forward diffusion process corrupts it with Gaussian noise to obtain $\mathbf{z}_t$ at timestep $t \in \{1,\dots,T\}$. 
A time-conditioned U-Net $\boldsymbol{\epsilon}_\theta$ is then trained to predict the noise $\boldsymbol{\epsilon}$ 
% \varun{i would change noise variable to $\eta$ to avoid confuson with the U-net} 
% \yian{kept $\epsilon$/$\epsilon_\theta$ for noise/U-Net to match the standard DDPM convention used in every concept-erasure baseline we compare against (ESD, STEREO, SAeUron, AdvUnlearn).}
added to $\mathbf{z}_t$. 
In text-conditioned generation, a prompt $p \in \mathcal{P}$ is encoded by a text encoder $f_{\psi}$ into text features $\mathbf{e}_p = f_{\psi}(p)$, and these features condition the U-Net through cross-attention layers. 
The standard latent diffusion objective is
\begin{equation}
\label{eq:ldm_loss}
\mathcal{L}_{\mathrm{LDM}}
=
\mathbb{E}_{x,\epsilon,t}
\left[
\left\|
\boldsymbol{\epsilon} - \boldsymbol{\epsilon}_\theta(\mathbf{z}_t, t, \mathbf{e}_p)
\right\|_2^2
\right],
\end{equation}
where $\epsilon \sim \mathcal{N}(0,I)$ 
%\ali{this is redundant: and $z_t$ is the noisy latent corresponding to $z = E(x)$}.
At inference time, generation starts from Gaussian noise and iteratively denoises toward a clean latent, which is finally decoded into an image.

% \ali{maybe "notation" in the name of this section is confusing because most of the notations and definitions will be in section 4.}

% \textbf{Contrast with vision-language models.}
% Although diffusion models use text encoders similar to those used in vision-language models (VLMs)~\citep{radford2021learning}, their training objectives differ fundamentally.
% VLMs such as CLIP are typically trained on large collections of image--text pairs using contrastive objectives that align image and text embeddings by pulling matched pairs together and pushing mismatched pairs apart~\citep{radford2021learning,zhang2024vision}. 
% Thus, VLM training primarily learns \emph{cross-modal alignment} in a shared representation space.
% In contrast, diffusion models use text embeddings as conditioning signals for a denoising generator: the text encoder does not directly predict image-text similarity, but instead provides token-level features that guide image synthesis through cross-attention.
% This distinction matters for unlearning: removing a concept from a generative diffusion model requires changing how text-conditioned denoising gives rise to visual content, not merely altering a retrieval-style similarity score.
% \textbf{Contrast with vision-language models.}
% \ali{I think this section is not necessary as we are not making any comparison or claim regarding the unlearning in VLMs.}
% Both vision-language models (VLMs)~\citep{radford2021learning} and text-to-image diffusion models learn a mapping between textual inputs and visual content. In VLMs, this mapping is used to align images and captions in a shared embedding space, while in diffusion models, text embeddings condition the denoising process that generates images. Thus, in both settings, text representations are tied to visual semantics, but diffusion models use this mapping generatively rather than for retrieval or matching. This distinction is important for unlearning: removing a concept from a diffusion model requires changing how text-conditioned representations give rise to visual content.


%\varun{i think the important part that is not highlighted is how there is text-image mapping in both models; the exact specifics of learning etc. are not relevant and can be removed/moved to appendix?}

% \textbf{Diffusion unlearning.}
% Let $\mathcal{D}_\theta$ denote a pretrained text-to-image diffusion model with parameters $\theta$, and let $c$ denote a target concept to be erased (e.g., \texttt{horse}).
% Diffusion unlearning seeks to produce an updated model $\mathcal{D}_{\theta'}$ such that prompts referring to $c$ no longer generate images containing that concept, while the model retains its ability to generate unrelated content.
% Formally, given a prompt distribution $\mathcal{P}_c$ associated with concept $c$, the goal is to reduce the probability that samples from $\mathcal{D}_{\theta'}(\cdot \mid p)$ contain $c$ for $p \sim \mathcal{P}_c$, while preserving generation quality and fidelity on prompts outside this target concept class. \varun{is there prior work on this? add some pointers, or maybe say this is discussed in intro? maybe this should be merged with the next subsection?}

\subsection{Unlearning Methods}
Machine unlearning in diffusion models aims to remove a target concept $c$ (e.g., \textit{horse}) from a pretrained model $\mathcal{D}_\theta$(the full text-to-image pipeline parameterized by the U-Net weights $\theta$ from Eq.~\ref{eq:ldm_loss}) 
% \varun{use a different variable for params, as the params $\theta$ are used for the U-Net?} 
% \yian{$\theta$ is shared with $\epsilon_\theta$ on purpose as the U-Net is the only trainable component, so $\mathcal{D}_\theta$ and $\epsilon_\theta$ must share parameters. }
while preserving overall generative capabilities. The goal is an updated model $\mathcal{D}_{\hat{\theta}}$ that suppresses $c$ on target-associated prompts while maintaining quality on unrelated prompts, following prior work on concept erasure in diffusion models. 

\textbf{Fine-tuning methods.}
A large body of work erase concepts by directly modifying the weights of a pretrained diffusion model. Erased Stable Diffusion (ESD)~\cite{gandikota2023erasing} uses negative guidance to align target-token prompts with a neutral anchor. Subsequent work improves the precision by editing cross-attention representations~\citep{wang2025ace} or updating salient concept-related parameters~\citep{fan2023salun,wu2024scissorhands}. More recent robust methods such as AdvUnlearn~\cite{zhang2024defensive} and STEREO~\cite{srivatsan2025stereo} further defend against prompt-based regeneration attacks.
However, these methods still largely define erasure through target tokens, prompts, or localized parameter subsets. They therefore suppress direct target prompts well, but struggle with \emph{concept leakage} through correlated context; when made more aggressive, their edits can propagate through shared representations and cause \emph{collateral damage} to semantic neighbors.

\textbf{Inference-time methods.}
A complementary line of work avoids modifying $\mathcal{D}_{\theta}$ and instead intervenes during generation by manipulating text embeddings~\citep{li2024get}, inserting lightweight concept-blocking adapters~\citep{lyu2024one}, or ablating sparse-autoencoder features~\citep{cywinski2025saeuron}. These methods are computationally efficient and often preserve broad model behavior, but their conservative, localized interventions do not remove the distributed contextual pathways through which a target concept can be recovered. As a result, they tend to preserve neighboring concepts better, while leaving residual indirect recovery.

\noindent{\bf Shared limitation.}
Thus, both fine-tuning and inference-time methods implicitly treat concepts as localized and independently suppressible. Our work challenges this assumption: when concepts occupy overlapping activation regions, erasure and retention are coupled, so stronger robustness against leakage necessarily increases the risk of neighbor degradation.

%\ali{we should somehow highlight this last paragraph; readers usually skip the related work sections.}
% \noindent\textbf{Takeaway.}
% Existing unlearning algorithms operate under the assumption that concepts can be selectively removed from the model’s representation space. In the following section, we provide empirical evidence that challenges this assumption by demonstrating that concepts remain recoverable through correlated contextual prompts.
% \noindent \textbf{Data Unlearning}
% \cite{alberti2025data}

\begin{comment}
    \subsection{Unlearning Methods}
% \noindent \textbf{The \textsc{UnlearnCanvas} Benchmark}
% Our experimental framework is built upon the \textsc{UnlearnCanvas} benchmark, a comprehensive dataset recently proposed by~\citet{zhang2024unlearncanvas} to standardize the evaluation of diffusion model unlearning. It was created to address the lack of rigorous, automated evaluation frameworks, which previously relied on ad-hoc, qualitative assessments. The benchmark consists of a high-resolution, stylized image dataset featuring 65 distinct artistic styles and 20 distinct object classes. Its primary advantage lies in its high stylistic consistency, which enables the training of a highly accurate style classifier, allowing for the reliable, quantitative measurement of unlearning using standardized metrics like Unlearning Accuracy (UA), In-domain Retain Accuracy (IRA), and Cross-domain Retain Accuracy (CRA). While \textsc{UnlearnCanvas} provides the essential visual data for our investigation, it does not include a formal, art style hierarchy linking its 60 styles. The styles are presented as a flat list. A foundational part of our methodology, therefore, is to construct this missing knowledge graph, which serves as the basis for our directional unlearning experiments.

\noindent \textbf{Methods for Diffusion Model Unlearning}
The primary objective of Machine Unlearning (MU) in diffusion models is to erase a target concept while minimizing collateral damage to the model's remaining generative capabilities. Current approaches can be broadly classified into two main categories:

\textbf{1. Fine-Tuning Methods:} This line of work involves permanently altering the model's weights. Many of these methods, such as Erased Stable Diffusion (ESD)~\cite{gandikota2023erasing}, utilize a form of negative guidance, retraining the model to align the output of the "forget" prompt with that of a neutral anchor prompt. Variations of this approach aim for greater precision. 
For example,~\citet{zhang2024forget} applies a re-steering loss to attention layers, while~\citet{fan2023salun} and~\citet{wu2024scissorhands} identify and modify only the most sensitive weights. 
\cite{zhang2024defensive} propose \textit{AdvUnlearn}, which integrates adversarial training into diffusion model unlearning to improve robustness against prompt-based attacks that attempt to regenerate erased concepts.
Attention-based Concept Erasure (ACE)~\cite{wang2025ace} removes target concepts by modifying cross-attention representations associated with the concept tokens, enabling localized editing of concept-related activations within the diffusion model.
STEREO~\cite{srivatsan2025stereo} introduces a robust concept unlearning framework that enforces separation between target and neighboring concepts, improving resistance to prompt-based regeneration of erased concepts.
%More recently,~\citet{li2025set} proposes a trajectory-aware finetuning framework that explicitly preserves early-stage denoising dynamics while reversing classifier-free guidance directions during mid-to-late stages, enabling precise concept erasure without relying on heuristic anchor selection.

Other notable fine-tuning methods include Unified Concept Editing (UCE)~\cite{gandikota2024unified}, Concept Ablation~\cite{vinker2023concept}, Selective Amnesia (SA)~\cite{heng2023selective}, and EraseDiff (EDiff)~\cite{wu2024erasediff}. 

While effective, these methods are vulnerable to "jailbreak" attacks that exploit related concepts~\cite{zhang2024generate} and operationally treat concepts as isolated entities. Orthogonal to concept-level unlearning,~\citet{alberti2025data} formalize datapoint-level unlearning in diffusion models and introduce Subtracted Importance Sampled Scores (SISS), an efficient fine-tuning objective with theoretical guarantees that aims to match the retrain-without-data gold standard at far lower cost~\cite{alberti2025data}

\textbf{2. Inference-Time Methods:} A second class of methods avoids the costly fine-tuning step by intervening directly during the generation process, leaving the model's original weights untouched. These techniques are often more efficient and can be dynamically toggled. SEOT~\cite{li2024get}suppresses unwanted content by manipulating text embeddings, and SPM~\cite{lyu2024one} uses lightweight one-dimensional adapters to block the propagation of concept-specific information. A highly relevant new approach in this category is SAeUron~\cite{cywinski2025saeuron}. This method leverages Sparse Autoencoders (SAEs) trained on the model's internal activations to learn a dictionary of sparse, interpretable features.And unlearning is achieved at inference time by identifying the features that correspond to an unwanted concept and ablating them, preventing their contribution to the final image. While these methods also treat concepts in isolation, their interventional nature inspires our work. We leverage similar interventions not just for deletion, but as a method to discover the latent structure of conceptual relationships. 
\end{comment}